\documentclass[11pt]{article}
\usepackage[T1]{fontenc}
\usepackage{lmodern}
\usepackage[margin=1in]{geometry}
\usepackage{amsmath,amssymb,amsthm,mathtools}
\usepackage{microtype,booktabs,tikz}
\usepackage[numbers,sort&compress]{natbib}
\usepackage{xurl}
\usepackage[colorlinks=true,linkcolor=blue!55!black,citecolor=blue!55!black,urlcolor=blue!55!black]{hyperref}
\hypersetup{pdftitle={Two-vertex certification of nonsingular sign-regular matrix intervals},pdfauthor={Wenqi Zhu; Kai Zou; Ping Nie}}
\newtheorem{theorem}{Theorem}[section]
\newtheorem{lemma}[theorem]{Lemma}
\newtheorem{proposition}[theorem]{Proposition}
\newtheorem{corollary}[theorem]{Corollary}
\theoremstyle{definition}
\newtheorem{example}[theorem]{Example}
\theoremstyle{remark}
\newtheorem{remark}[theorem]{Remark}
\numberwithin{equation}{section}
\newcommand{\R}{\mathbb R}
\newcommand{\eps}{\varepsilon}
\newcommand{\K}{\mathcal K}
\newcommand{\chle}{\leq_{\mathrm c}}
\newcommand{\one}{\mathbf 1}
\DeclareMathOperator{\rank}{rank}
\DeclareMathOperator{\diag}{diag}
\DeclareMathOperator{\ri}{ri}
\newcommand{\cornerdiagram}[3]{%
  \begin{tikzpicture}[x=.58cm,y=-.58cm,font=\scriptsize]
  \fill[blue!12] #2 rectangle #3;
  \draw (0,0) rectangle (4,4);
  \draw[dashed,gray] (1.2,1.2) rectangle (2.8,2.8);
  \foreach \x in {1.2,2.8}{\foreach \y in {1.2,2.8}{\fill (\x,\y) circle (2pt);}}
  \node at (1.2,-.35) {$k$}; \node at (2.8,-.35) {$\ell$};
  \node at (-.35,1.2) {$i$}; \node at (-.35,2.8) {$j$};
  \node at (2,4.55) {#1};
  \end{tikzpicture}}
\title{Two-vertex certification of nonsingular\ sign-regular matrix intervals}
\author{%
Wenqi Zhu\\
{\small Mathematical Institute, Woodstock Road, University of Oxford,}\\
{\small Oxford, UK, OX2 6GG}\\
{\small\texttt{wenqi.zhu@maths.ox.ac.uk}}\\[1ex]
Kai Zou\\
{\small Mathematical \& Engineering Sciences, King's College London, UK}\\
{\small\texttt{kz@netmind.ai}}\\[1ex]
Ping Nie\\
{\small David R. Cheriton School of Computer Science, University of Waterloo}\\
{\small\texttt{ping.nie@uwaterloo.ca}}}
\date{}

\begin{document}
\maketitle
\begin{abstract}
For every positive integer $n$ and every signature
$\eps\in\{-1,1\}^n$, we prove that a real square entrywise matrix box
consists entirely of nonsingular sign-regular matrices with signature
$\eps$ if and only if its two checker vertices have that property.
Zero minors of every proper order and fixed entries of both checker
parities are allowed. This proves the general interval conjecture
stated by Garloff, Adm, and Titi in 2016. The proof uses induction on
minor order. A first wrong-sign minor of minimal span forces a
rank-deficient rectangle whose rank bound persists throughout the box.
Independent entry variation then produces an essential row or column.
Signed Schur complements and rank-one corner structure exclude these
variations; a fixed rectangle is excluded by a common corner and an
affine midpoint argument.
\end{abstract}

\section{The main theorem and its context}\label{sec:intro}

For increasing index sets $I,J$ of equal cardinality, write $X[I,J]$
for the corresponding submatrix, with its rows and columns in increasing
order. A real $n\times n$ matrix $X$ is \emph{sign regular with
signature} $\eps=(\eps_1,\ldots,\eps_n)\in\{-1,1\}^n$ if
\begin{equation}\label{eq:sr}
  \eps_k\det X[I,J]\geq0
  \qquad (|I|=|J|=k,\ 1\leq k\leq n).
\end{equation}
Thus zero minors are allowed. Strict sign regularity requires strict
inequalities in~\eqref{eq:sr}. Nonsingular sign regularity requires
\eqref{eq:sr} and $\det X\neq0$.

Let $L,U\in\R^{n\times n}$ satisfy $L\leq U$ entrywise. Define the
\emph{checker vertices} of the entrywise box $[L,U]$ by
\begin{equation}\label{eq:vertices}
 A^-_{ij}=\begin{cases}L_{ij},&i+j\text{ even},\\U_{ij},&i+j\text{ odd},\end{cases}
 \qquad
 A^+_{ij}=\begin{cases}U_{ij},&i+j\text{ even},\\L_{ij},&i+j\text{ odd}.
 \end{cases}
\end{equation}

\begin{theorem}\label{thm:main}
Let $n\geq1$, let $\eps\in\{-1,1\}^n$, and let $L,U\in\R^{n\times n}$
satisfy $L\leq U$. The following statements are equivalent:
\begin{enumerate}
\item[(i)] Every $X\in[L,U]$ is nonsingular and sign regular with signature $\eps$.
\item[(ii)] Both matrices $A^-$ and $A^+$ in~\eqref{eq:vertices} are nonsingular
and sign regular with signature $\eps$.
\end{enumerate}
\end{theorem}

The implication from (i) to (ii) follows because the two vertices belong
to the box. The content of the theorem is the converse, with no restriction
on the positions of fixed entries or vanishing proper minors. In
particular, a fixed entry need not have the same checker parity as every
other fixed entry.

\subsection{Earlier interval results}

The two-vertex interval property for strictly sign-regular matrices is
classical. Garloff~\cite{Garloff1982} established this property and a
weak-sign version under a restriction on fixed positions: all fixed
entries must have even index sum, or all must have odd index sum.
The latter includes boxes without fixed entries; see also
\cite[Theorem~3.1]{Survey2016} and
\cite[Theorems~3.2--3.3]{AdmGarloff2026}.

Garloff's 1982 conjecture for nonsingular totally nonnegative matrices
was proved by Adm and Garloff~\cite[Theorem~3.6]{AdmGarloff2013}.
Here total nonnegativity means that every minor is nonnegative.
Their proof uses the Cauchon algorithm, and their Theorem~3.4 also
preserves a common zero/nonzero pattern of all minors. Earlier special
cases included almost totally positive
matrices~\cite{Garloff2003}. Related singular cases require additional
conditions: see \cite[Corollaries~3.7--3.8]{AdmGarloff2013} and the
rank and independence hypotheses in~\cite{Relax2020}.

Adm and Garloff~\cite{AdmGarloff2016} proved the interval property for
nonsingular almost strictly sign-regular matrices and for nonsingular
tridiagonal sign-regular matrices of a common signature. They also treated
totally nonpositive matrices and related signatures. Further signatures,
including those of matrices with nonnegative proper minors and negative
determinant, were treated in~\cite{AdmGarloff2021}; see
\cite[IP~2.8--2.9]{Survey2021}. Hassuneh, Adm, and
Garloff~\cite{HassunehAdmGarloff2025} considered matrices with
nonpositive proper minors and positive determinant. Some of those
results retain a strict condition on a corner entry. The resulting
lists involve eight period-four signatures and their final-sign
reversals. The survey~\cite{AdmGarloff2026} records the coverage of
small orders together with the relevant strict-entry qualifications.

The general statement is Conjecture~3.1 of Garloff, Adm, and
Titi~\cite[p.~7]{Survey2016}: nonsingular sign-regular matrices with a
fixed signature should possess the interval property. This general
conjecture is distinct from the totally nonnegative conjecture already
settled in~\cite{AdmGarloff2013}.

There are general tests using more vertices. For an $m\times n$ box,
write
\[
 A_c=\tfrac12(L+U),\qquad \Delta=\tfrac12(U-L).
\]
Adm and Garloff~\cite[Theorem~3.1]{AdmGarloff2025} proved that the box
is sign regular with a prescribed signature if and only if all matrices
\begin{equation}\label{eq:manyvertices}
 A_c-\diag(y)\Delta\diag(z),
 \qquad y\in\{-1,1\}^m,\quad z\in\{-1,1\}^n,
\end{equation}
have that signature. There are at most $2^{m+n-1}$ distinct matrices
in~\eqref{eq:manyvertices}. Their argument uses variation diminution;
for related single-vector characterizations and interval-hull results,
see~\cite{ChoudhuryKannanKhare2021,ChoudhuryYadav2025}.
The January 2026 survey~\cite{AdmGarloff2026} explicitly leaves open
whether, in the square nonsingular case, the larger vertex family can
be reduced to the two checker vertices. Theorem~\ref{thm:main} gives
that reduction.

Two sufficient conditions in the public September 2026 manuscripts of
Holden are closer to the boundary cases addressed here.
Theorem~G of~\cite{HoldenGraph2026} assumes an acyclic directed graph
of nonzero fixed entries and confines fixed zero entries to one checker
parity. A positive diagonal scaling then reduces to the older
fixed-parity result. Theorem~V of~\cite{HoldenFive2026} assumes no fixed
zero entries and requires every fully fixed adjacent $2\times2$ block
to be nonsingular. It controls order-three minors in every dimension
and proves the full conclusion at order five. The unconditional
order-two reduction in~\cite{HoldenTwo2026}, together with the
inverse-positive reduction in~\cite{HoldenGraph2026}, supplies useful
starting points. We include proofs of both reductions below. The higher
orders are treated without the additional hypotheses in those manuscripts.

\subsection{The boundary obstruction}

The main issue is compatibility with the box. A nonsingular
sign-regular matrix can be approximated by strictly sign-regular
matrices with the same signature. Approximating the two checker
vertices separately, however, does not ensure that the approximants
remain checker-comparable in a way that yields the original interval
conclusion. This difficulty is described explicitly
in~\cite{AdmGarloff2026}.

The proof below treats zero minors directly. Its determinant identities,
inverse-positive arguments, and rank factorizations are classical.
The corner structure is related to the shadow property for totally
nonnegative matrices; see~\cite[Lemma~1.2]{AdmGarloff2013} and the
background in~\cite{FallatJohnson2011,Pinkus2010}. We use these ideas
in a form that depends only on the low-order signs visible after Schur
elimination and that handles both signs of the resulting third-order
minors.

\section{Notation and proof strategy}\label{sec:notation}

All matrix indices start at $1$. Index sets are always written in
increasing order when a determinant is taken. For $a\leq b$, write
$[a,b]=\{a,a+1,\ldots,b\}$ when an integer index set is intended, and
$[n]=[1,n]$. Complements of index sets are taken in the full row or
column index set of the matrix under discussion. We put $\eps_0=1$
and give the empty determinant the value $1$.

Let
\[
 D_n=\diag((-1)^1,\ldots,(-1)^n).
\]
The \emph{checker order} is defined by
\[
 A\chle B\quad\Longleftrightarrow\quad D_nAD_n\leq D_nBD_n.
\]
For $A\chle B$, put
\[
 \K(A,B)=\{X:A\chle X\chle B\}.
\]
This is an entrywise Cartesian box, whose ordinary lower and upper
bounds need not be $A$ and $B$. For the box in
Theorem~\ref{thm:main}, they are the checker vertices $A=A^-$ and
$B=A^+$. A coordinate is \emph{fixed} if $A_{ij}=B_{ij}$ and
\emph{free} otherwise. A matrix lies in $\ri\K$ if each free
coordinate lies strictly between its two bounds. Thus relative interior
is taken in the affine space of the actual box, including when some or
all coordinates are fixed.

A \emph{solid} minor uses consecutive rows and consecutive columns.
An \emph{adjacent two-minor} is a solid minor of order two.
We use $\mathrm{TN}_2$ to mean a nonnegative matrix whose two-minors
are all nonnegative; it carries no strictness assertion.

Here is the structure of the proof. Section~\ref{sec:initial} proves
nonsingularity throughout the box and the prescribed signs through
order two. In particular, it also gives the cofactor signs of order
$n-1$. Sections~\ref{sec:corners} and~\ref{sec:schur} establish the
low-order corner and Schur-complement facts used in the induction.
Section~\ref{sec:rank} shows that a wrong-sign minor of minimal span
forces a permanent rank bound on an interior-column rectangle. If the
rectangle has a free entry, an essential row or column appears.
Section~\ref{sec:induction} rules out such a line by Schur elimination,
and treats a fixed rectangle using a common rank-one corner and the
midpoint of the box. Examples and scope remarks conclude the paper.

The reductions already control orders $1,2,n-1,n$. These exhaust the
minor orders for $n\leq4$; at $n=5$, order three is the first remaining
order. The induction treats all subsequent orders by the same argument.

\section{Initial interval facts}\label{sec:initial}

\subsection{Nonsingularity and cofactor signs}

\begin{lemma}\label{lem:inverse}
Suppose $P,Q\in\R^{n\times n}$ are nonsingular, $P\leq Q$, and
$P^{-1},Q^{-1}\geq0$. Every $C$ with $P\leq C\leq Q$ is
nonsingular and satisfies $C^{-1}\geq0$.
\end{lemma}

\begin{proof}
Set
\[
 v=P^{-1}\one,\qquad w=Q^{-1}\one,\qquad T=Q^{-1}(Q-P).
\]
Every row of an invertible nonnegative matrix has a positive entry,
so $v,w$ are strictly positive. Moreover, $T\geq0$ and
$Tv=v-w<v$. In the weighted maximum norm
$\|x\|_v=\max_i|x_i|/v_i$, the norm of a nonnegative matrix satisfies
\[
 \|T\|_v=\max_i\frac{(Tv)_i}{v_i}<1.
\]
For $P\leq C\leq Q$, let $E=Q^{-1}(Q-C)$. Then
$0\leq E\leq T$, and hence $\|E\|_v<1$. Since $C=Q(I-E)$,
the Neumann series yields
\[
 C^{-1}=\left(\sum_{s=0}^{\infty}E^s\right)Q^{-1}\geq0.
\]
\end{proof}

\begin{corollary}\label{cor:nonsingular}
Suppose $A\chle B$ are nonsingular sign-regular matrices with
signature $\eps$. Every $X\in\K(A,B)$ is nonsingular, its determinant
has sign $\eps_n$, and its $(n-1)$-minors have their prescribed weak
sign $\eps_{n-1}$.
\end{corollary}

\begin{proof}
The assertion about order zero is interpreted by the empty determinant
when $n=1$. Put $\delta=\eps_{n-1}\eps_n$. The inverse cofactor formula
gives
\[
 \delta D_nA^{-1}D_n\geq0,\qquad
 \delta D_nB^{-1}D_n\geq0.
\]
If $\delta=1$, apply Lemma~\ref{lem:inverse} between $D_nAD_n$ and
$D_nBD_n$. If $\delta=-1$, apply it between $-D_nBD_n$ and
$-D_nAD_n$. Thus every $X\in\K$ is nonsingular and satisfies
$\delta D_nX^{-1}D_n\geq0$. Since $\K$ is connected, its nowhere-zero
determinant has the sign of $\det A$, namely $\eps_n$.
Applying the inverse cofactor formula once more gives the signs of the
$(n-1)$-minors.
\end{proof}

\subsection{Staircase support}

\begin{lemma}\label{lem:staircase}
Let $Y\in\R^{N\times N}$ be nonsingular and $\mathrm{TN}_2$.
Every diagonal entry of $Y$ is positive. A zero $Y_{ij}=0$ with $i<j$
forces $Y_{rc}=0$ for $r\leq i$ and $c\geq j$. A zero below the
diagonal forces all entries southwest of it to vanish. Consequently,
the positive entries of every row and column form a nonempty integer
interval containing its diagonal position. The left and right endpoints
of the row intervals are nondecreasing with the row index.
\end{lemma}

\begin{proof}
Nonsingularity supplies a permutation $\pi$ such that
$\prod_iY_{i,\pi(i)}>0$. If $\pi$ has an adjacent inversion,
$\pi(i)>\pi(i+1)$, the two-minor on these rows and columns gives
\[
 Y_{i,\pi(i+1)}Y_{i+1,\pi(i)}
 \geq Y_{i,\pi(i)}Y_{i+1,\pi(i+1)}>0.
\]
Swapping these two permutation values therefore preserves positivity
of the selected product. Repeatedly removing adjacent inversions
reaches the identity permutation, proving $Y_{ii}>0$ for every $i$.

Suppose $Y_{ij}=0$ with $i<j$. For $c>j$, the two-minor on rows
$i,j$ and columns $j,c$ reads
$-Y_{ic}Y_{jj}\geq0$, so $Y_{ic}=0$; the same conclusion for
$c=j$ is the assumption. If $r<i$ and $c\geq j$, the minor on
rows $r,i$ and columns $i,c$ now reads
$-Y_{rc}Y_{ii}\geq0$. This proves northeast zero propagation.
Transposition proves the southwest assertion.

A row cannot have a zero between two of its positive entries: to the
right of its positive diagonal such a zero would force the rest of the
row to vanish, and to the left it would force all earlier entries to
vanish. Thus its support is an interval containing the diagonal.
If a lower row began to the left of an earlier row, the southwest
propagation from a zero just before the earlier row's support would
give a contradiction. If an earlier row ended to the right of a later
row, northeast propagation from the later row gives the analogous
contradiction. Hence both support endpoints are nondecreasing.
The column assertion follows by transposition.
\end{proof}

We refer to these support properties as \emph{staircase support}.
They will also be applied after Schur elimination and a sign
normalization.

\subsection{The first two minor orders}

\begin{proposition}\label{prop:ordertwo}
Under the hypotheses of Corollary~\ref{cor:nonsingular}, every matrix
in $\K(A,B)$ has the prescribed signs at orders one and two.
\end{proposition}

\begin{proof}
Order one follows entry by entry. To treat order two, negate all
matrices if necessary so that their entries are nonnegative. Reversing
all rows, if necessary, then makes the endpoint two-minors
nonnegative. Negation multiplies every $k$-minor by $(-1)^k$, and
row reversal multiplies it by $(-1)^{k(k-1)/2}$. Both operations map a
checker box to a checker box, possibly exchanging its two vertices.
We may therefore work with nonnegative endpoints whose two-minors
are nonnegative.

On consecutive rows and columns, the determinant $ad-bc$ is
nondecreasing in $a,d$ and nonincreasing in $b,c$. Its two extrema
over this block of the box occur at the restrictions of the two checker
vertices. All adjacent two-minors throughout the box are therefore
nonnegative.

Take $X\in\ri\K$. Its positive support is the union of the endpoint
supports. Each row's two endpoint support intervals contain its
diagonal position. Their union is thus an interval; its left endpoint
is the minimum of two nondecreasing functions of the row, and its
right endpoint is their maximum. Hence this union also has staircase
support.

Consider rows $i<j$ and columns $k<\ell$. If either $X_{i\ell}$ or
$X_{jk}$ is zero, the desired inequality
\begin{equation}\label{eq:two}
 X_{ik}X_{j\ell}\geq X_{i\ell}X_{jk}
\end{equation}
is immediate. Otherwise every entry of $X$ in
$[i,j]\times[k,\ell]$ is positive. Indeed, for each intermediate row,
its left support endpoint is at most that of row $j$, hence at most
$k$, and its right endpoint is at least that of row $i$, hence at
least $\ell$. Multiply the adjacent inequalities
\[
 X_{rc}X_{r+1,c+1}\geq X_{r,c+1}X_{r+1,c},
 \qquad i\leq r<j,\quad k\leq c<\ell.
\]
All factors that occur on both sides are positive and can be
cancelled. The resulting inequality is~\eqref{eq:two}.
Continuity extends it from $\ri\K$ to $\K$. Undoing the
normalizations proves the stated signs.
\end{proof}

\section{Low-order corner structure}\label{sec:corners}

In this section $Y\in\R^{N\times N}$ is nonsingular and
$\mathrm{TN}_2$, with $N\geq3$. When a sign $\eta\in\{-1,1\}$
is specified, all three-minors of $Y$ have weak sign $\eta$.

\subsection{A positive vanishing two-minor}

\begin{lemma}\label{lem:corner}
Suppose $i<j$, $k<\ell$, all four entries of
$Y[\{i,j\},\{k,\ell\}]$ are positive, and its determinant is zero.
If the three-minors have weak sign $\eta$, one of the following
rectangles has rank one:
\begin{align*}
 \eta=-1:&\quad [1,j]\times[1,\ell]\quad\text{or}\quad[i,N]\times[k,N],\\
 \eta=+1:&\quad [1,j]\times[k,N]\quad\text{or}\quad[i,N]\times[1,\ell].
\end{align*}
For these fixed indices and this value of $\eta$, at most one of the
two displayed orientations can occur in a nonsingular matrix.
\end{lemma}

\begin{figure}[htbp]
\centering
\begin{tabular}{cccc}
\cornerdiagram{$\eta=-1$: NW}{(0,0)}{(2.8,2.8)}&
\cornerdiagram{$\eta=-1$: SE}{(1.2,1.2)}{(4,4)}&
\cornerdiagram{$\eta=+1$: NE}{(1.2,0)}{(4,2.8)}&
\cornerdiagram{$\eta=+1$: SW}{(0,1.2)}{(2.8,4)}
\end{tabular}
\caption{The four possible rank-one corner orientations. The dots
mark the positive entries of the vanishing two-minor. Rows increase
downward and columns to the right. These are schematic orientations;
the rank bound in the proof excludes one member of the appropriate
pair for any fixed set of indices.}\label{fig:corners}
\end{figure}

\begin{proof}
Divide rows $i$ and $j$ by their positive column-$k$ entries, and
denote the normalized rows by $u$ and $v$. There is a number $\rho>0$
such that
\[
 u_k=v_k=1,\qquad u_\ell=v_\ell=\rho.
\]
Put $d_t=v_t-u_t$. The two-minors of these rows give
\begin{equation}\label{eq:d-signs}
 d_t\leq0\ (t<k),\qquad
 d_t=0\ (k\leq t\leq\ell),\qquad
 d_t\geq0\ (t>\ell).
\end{equation}
For example, inside $[k,\ell]$ the minors using column $k$ give
$d_t\geq0$ and those using column $\ell$ give $d_t\leq0$.
For each row $h$, put
\[
 q_h=Y_{h\ell}-\rho Y_{hk}.
\]
Using the two-minors of columns $k,\ell$ and either anchor row gives
\begin{equation}\label{eq:q-signs}
 q_h\leq0\ (h<i),\qquad
 q_h=0\ (i\leq h\leq j),\qquad
 q_h\geq0\ (h>j).
\end{equation}
Some $q_h$ is nonzero, because columns $k,\ell$ of the nonsingular
matrix $Y$ are linearly independent.

For $h<i$ or $h>j$ and $t<k$ or $t>\ell$, consider the three-minor
on these three rows and columns, sorted in increasing order. After
the two positive row normalizations, subtracting $u$ from $v$ shows
that its value is
\begin{equation}\label{eq:dq}
 -d_tq_h.
\end{equation}
For instance, in rows $i,j,h$ and columns $k,\ell,t$, the second row
becomes $(0,0,d_t)$, whose cofactor is $-q_h$.
The other three placements give the same sign. The identity does not
require entries of row $h$ to be positive.

If both an earlier row with $q_h<0$ and a later row with $q_h>0$
existed, the common weak sign $\eta$ in~\eqref{eq:dq} would force
$d_t=0$ on both outer column ranges. Together with
\eqref{eq:d-signs}, this would make the two full rows proportional,
contradicting nonsingularity. Thus the strict $q_h$ occur on exactly
one side of the anchor rows.

Suppose first that a later row $h>j$ has $q_h>0$. Then $q_r=0$
for every $r\leq j$. Consider such a row with $Y_{rk}>0$.
Its normalized entries in columns $k,\ell$ are again $1,\rho$.
For $r<j$, apply~\eqref{eq:dq} to the pair of rows $r,j$ and
the later row $h$. If $\eta=+1$, the identity forces the normalized
row difference to be nonpositive, while the two-minors make that
difference nonnegative on columns $t>\ell$. It therefore vanishes
there and throughout $[k,\ell]$. Thus row $r$ is proportional to
row $j$ on all columns $t\geq k$. If $\eta=-1$, the same identity
forces a nonnegative difference, which must vanish for $t<k$ and
throughout $[k,\ell]$. In this case the rows are proportional on
all columns $t\leq\ell$. The conclusion for $r=j$ is immediate.

It remains to handle $r\leq j$ with $Y_{rk}=0$. Equation
\eqref{eq:q-signs} gives $Y_{r\ell}=0$. Staircase support and
$Y_{jk}>0$ imply $r<k$: a zero at or below the diagonal in column
$k$ could not have a positive entry farther down that column.
Also $r<i$, since the column support contains every row between the
two positive anchor entries. Northeast zero propagation now shows
that row $r$ is zero in every column $t\geq k$. When $\eta=+1$,
this is precisely a zero row of the asserted northeast corner.
When $\eta=-1$, it is impossible: on rows $r,i,h$ and columns
$r,k,\ell$, the determinant is
\[
 Y_{rr}\det Y[\{i,h\},\{k,\ell\}]
 =Y_{rr}Y_{ik}q_h>0.
\]
This contradicts the nonpositive third-order sign.
Consequently the northeast corner has rank one for $\eta=+1$, and
the northwest corner has rank one for $\eta=-1$. Its rank is exactly
one because it contains the positive anchor entries.

If instead the strict $q_h$ occur on the earlier side, reverse both
rows and columns. A minor of order $s$ acquires the factor
$(-1)^{s(s-1)}=1$, so its sign is preserved. The preceding argument
then gives the southwest or southeast corner, respectively.

For uniqueness, any $p\times q$ submatrix of an invertible
$N\times N$ matrix has rank at least $p+q-N$: adjoining the missing
rows and columns can increase rank by at most $(N-p)+(N-q)$.
A rank-one corner therefore has $p+q\leq N+1$. For $\eta=-1$,
a northwest corner would require $j+\ell\leq N+1$, while a southeast
corner would require $i+k\geq N+1$. These inequalities cannot both
hold because $i+k<j+\ell$. For $\eta=+1$, northeast and southwest
would require $j\leq k$ and $\ell\leq i$, respectively, which are
also incompatible. These necessary inequalities depend only on the
anchor indices. They exclude different orientations even in two
different nonsingular matrices with the same anchor indices and
third-order sign.
\end{proof}

\subsection{Zeros when the third-order sign is negative}

\begin{lemma}\label{lem:twozeros}
If all three-minors of $Y$ are nonpositive, its only possible zero
entries are $(1,N)$ and $(N,1)$.
\end{lemma}

\begin{proof}
The diagonal is positive by Lemma~\ref{lem:staircase}.
Suppose $Y_{ij}=0$ with $i<j$. The rectangle
$[1,i]\times[j,N]$ is zero. If $N-j+1\geq2$, select two columns
with indices at least $j$. They are independent because $Y$ is
nonsingular. They both vanish on the first $i$ rows, so some
two-minor on rows greater than $i$ is nonzero, hence positive.
Append row $i$ and column $i$ to this minor. The resulting sorted
three-minor is $Y_{ii}$ times the positive two-minor, contradicting
the assumed sign. Hence $j=N$.

Similarly, if $i\geq2$, two rows with indices at most $i$ are
independent and both vanish on columns $j,\ldots,N$. Some two-minor
in columns less than $j$ is therefore positive. Appending row $j$
and column $j$ gives another positive three-minor. Thus $i=1$.
Transposition treats zeros below the diagonal.
\end{proof}

\begin{lemma}\label{lem:middle}
A $3\times3$ $\mathrm{TN}_2$ matrix having a zero in its middle
column has nonnegative determinant.
\end{lemma}

\begin{proof}
Write the matrix as
$\left(\begin{smallmatrix}a&b&c\\d&e&f\\g&h&i\end{smallmatrix}\right)$.
If $e=0$, two-minor nonnegativity gives $bd=fh=0$. Expanding the
determinant leaves $bfg+cdh\geq0$. If $e>0$ and $b=0$, the
upper-right two-minor gives $c=0$, and the determinant is
$a(ei-fh)\geq0$. If $e>0$ and $h=0$, the lower-left two-minor
gives $g=0$, and the determinant is $i(ae-bd)\geq0$. These cases
cover every zero in the middle column.
\end{proof}

\section{Signed Schur complements and rank completion}\label{sec:schur}

\subsection{The ordering signs in Schur elimination}

Let $Z\in\R^{n\times n}$ be nonsingular. Choose row and column
sets $P,Q$ with $|P|=|Q|=p$ such that $H=Z[P,Q]$ is
nonsingular. The pivot need not be principal or contiguous.
In the complementary index orders inherited from $Z$, set
\begin{equation}\label{eq:schur}
 S_Z=Z[P^c,Q^c]-Z[P^c,Q]H^{-1}Z[P,Q^c].
\end{equation}
For $p=0$, the pivot is empty, its determinant is $1$, and we
interpret $S_Z=Z$. All subsequent formulas include this case.

For $i\notin P$ and $j\notin Q$, define
\begin{equation}\label{eq:ordering-signs}
 \alpha_i=(-1)^{|\{u\in P:u>i\}|},\qquad
 \beta_j=(-1)^{|\{v\in Q:v>j\}|}.
\end{equation}
Let $\widehat S_Z=\diag(\alpha)S_Z\diag(\beta)$.

\begin{lemma}\label{lem:schur}
For increasing sets $I\subseteq P^c$, $J\subseteq Q^c$ of equal
cardinality,
\begin{equation}\label{eq:bordered}
 \det\widehat S_Z[I,J]
   =\frac{\det Z[P\cup I,Q\cup J]}{\det H}.
\end{equation}
For any $U\supseteq P$ and $V\supseteq Q$, with no requirement
that $|U|=|V|$,
\begin{equation}\label{eq:rank-lifting}
 \rank Z[U,V]
   =p+\rank S_Z[U\setminus P,V\setminus Q].
\end{equation}
In particular, $S_Z$ is nonsingular.
\end{lemma}

\begin{proof}
First order the bordered rows by listing $P$ followed by $I$, each
internally increasing, and do the same with $Q$ followed by $J$.
Block elimination then gives the determinant
$\det H\det S_Z[I,J]$. Sorting the combined row list has sign
$\prod_{i\in I}\alpha_i$: the entry $i$ crosses exactly the pivot
rows larger than $i$. The column sorting sign is
$\prod_{j\in J}\beta_j$. Thus
\[
 \det H\det S_Z[I,J]
 =\left(\prod_{i\in I}\alpha_i\right)
  \left(\prod_{j\in J}\beta_j\right)
  \det Z[P\cup I,Q\cup J].
\]
Multiplying by these two sign products gives~\eqref{eq:bordered}.

For~\eqref{eq:rank-lifting}, perform the same row and column
permutations and invertible block elimination on $Z[U,V]$.
The result is block diagonal with blocks $H$ and
$S_Z[U\setminus P,V\setminus Q]$. Their ranks add. Taking $U=V=[n]$
proves nonsingularity of $S_Z$.
\end{proof}

Suppose now that $\det H$ has sign $\eps_p$, and that the relevant
minors of orders $p+1$ and $p+2$ have their prescribed signs.
Define
\begin{equation}\label{eq:normalized-schur}
 T_Z=(\eps_p\eps_{p+1})\widehat S_Z.
\end{equation}
Its entries are nonnegative. Equation~\eqref{eq:bordered} gives its
order-two sign as
\[
 \tau_2=\eps_p\eps_{p+2}.
\]
If the order-$(p+3)$ signs are prescribed as well, its order-three
sign is
\[
 \tau_3=\eps_{p+1}\eps_{p+3}.
\]
Indeed the scalar factor in~\eqref{eq:normalized-schur} is raised
to the minor order, and the denominator in~\eqref{eq:bordered}
has sign $\eps_p$.

If $\tau_2=-1$, reverse all rows of $T_Z$. Row reversal multiplies
both two-minors and three-minors by $-1$ and preserves entry signs.
If $\tau_2=+1$, leave the rows in place. We denote the resulting
matrix by $\mathcal T_Z$. It is nonsingular and $\mathrm{TN}_2$.
Its prescribed third-order sign, when that order is known, is
\begin{equation}\label{eq:delta}
 \delta_p=\eps_p\eps_{p+1}\eps_{p+2}\eps_{p+3}.
\end{equation}
The sign bookkeeping can be summarized as follows:
\begin{center}
\begin{tabular}{ccl}
\toprule
minor order & in $T_Z$ & in $\mathcal T_Z$\\
\midrule
$1$ & $+1$ & $+1$\\
$2$ & $\eps_p\eps_{p+2}$ & $+1$\\
$3$ & $\eps_{p+1}\eps_{p+3}$ & $\delta_p$\\
\bottomrule
\end{tabular}
\end{center}
After reversal, row indices and increasing row order refer to the
new row positions; column order is unchanged.

Only the already established signs at orders $p+1$ and $p+2$ are
needed to conclude that $\mathcal T_Z$ is $\mathrm{TN}_2$.
If an interior matrix has a wrong-sign bordered $(p+3)$-minor,
the corresponding three-minor of $\mathcal T_Z$ has the wrong sign
relative to~\eqref{eq:delta}. No third-order sign condition on this
interior Schur matrix is being assumed. Further, the Schur matrices
of checker-comparable matrices need not themselves be
checker-comparable; the proof uses no such assertion.

\subsection{Completion of a deficient pivot}

\begin{lemma}\label{lem:rank-obstruction}
Let $Z\in\R^{n\times n}$ be nonsingular and sign regular with
signature $\eps$. Let $0\leq p\leq n-3$ and suppose
$\delta_p=-1$ in~\eqref{eq:delta}. If
\[
 \rank Z[R_0,C_0]=t\leq p,\qquad |R_0|\geq p,
\]
then
\begin{equation}\label{eq:zero-area}
 (|R_0|-p)(|C_0|-t)\leq2.
\end{equation}
In particular, a rectangle of rank at most $p$ cannot have either
of the following sizes:
\begin{equation}\label{eq:forbidden-sizes}
 \text{at least }(p+2)\times(p+2),\qquad
 \text{at least }(p+3)\times(p+1).
\end{equation}
\end{lemma}

\begin{proof}
Choose $t$ columns in $C_0$ that form a basis of the columns
restricted to $R_0$. The full rows $R_0$ of the nonsingular matrix
$Z$ are independent. Hence the matrix $Z[R_0,[n]]$ has rank
$|R_0|\geq p$, and the selected $t$ columns can be extended to
$p$ independent columns on $R_0$. Denote their index set by $L$.
Every added column can be chosen outside $C_0$, because every
column in $C_0$ already lies in the span of the initial $t$ columns.
Consequently,
\[
 |L|=p,\qquad |L\cap C_0|=t.
\]
Choose $p$ rows $P_0\subseteq R_0$ for which $Z[P_0,L]$ is
invertible. This is the pivot used for Schur elimination; it is
formed by completing the basis when $t<p$, rather than assuming
that the original rectangle has rank $p$.

For each column $c\in C_0$, its entries on all of $R_0$ are a
linear combination of the initial $t$ basis columns. Extending that
coefficient vector by zeros in the other $p-t$ positions of $L$
and restricting to $P_0$ shows that these are exactly the coefficients
computed by $Z[P_0,L]^{-1}Z[P_0,c]$. Thus its Schur residual is
zero on $R_0\setminus P_0$. The Schur matrix has a zero rectangle
on
\[
 (R_0\setminus P_0)\times(C_0\setminus L)
\]
of size $(|R_0|-p)\times(|C_0|-t)$.

The normalized Schur matrix has order $n-p\geq3$, is nonsingular
and $\mathrm{TN}_2$, and has nonpositive three-minors. Diagonal
sign changes and row reversal preserve zero entries. By
Lemma~\ref{lem:twozeros}, this whole Schur matrix has at most two
zero entries. This proves~\eqref{eq:zero-area}.
If the first size in~\eqref{eq:forbidden-sizes} occurred, the zero
rectangle would have at least $2\times2$ entries. For the second
size it would have at least $3\times1$ entries. Both are impossible.

When $p=0$, necessarily $t=0$; the selected column and row sets
are empty and no elimination is needed. The rectangle itself is
zero in the normalized matrix, and exactly the same argument
applies. This also verifies the empty-pivot case directly.
\end{proof}

\section{From a bad minor to a permanent rank bound}\label{sec:rank}

Fix $k\geq3$. For the remainder of this section, suppose that the
endpoint hypotheses hold and that every minor of order at most
$k-1$ has its prescribed weak sign throughout $\K=\K(A,B)$.

\subsection{Interior zeros of multiaffine minors}

\begin{lemma}\label{lem:permanent}
If a minor of order at most $k-1$ vanishes at a point of $\ri\K$,
it vanishes throughout $\K$. Consequently, if a submatrix has rank
$s\leq k-2$ at a relative-interior point, its rank is at most $s$
at every point of $\K$.
\end{lemma}

\begin{proof}
List the genuinely free scalar coordinates of $\K$ as
$x_1,\ldots,x_d$, each varying in a nondegenerate interval
$[a_q,b_q]$. At an interior coordinate let
$\theta_q=(x_q-a_q)/(b_q-a_q)\in(0,1)$. A determinant is affine
in each scalar entry separately. Repeated interpolation therefore
expresses its signed value $f$ as
\begin{equation}\label{eq:interpolate}
 f(x)=\sum_{\omega\in\{0,1\}^d}
 \left(\prod_{q=1}^d\theta_q^{\omega_q}
            (1-\theta_q)^{1-\omega_q}\right) f(v_\omega),
\end{equation}
where $v_\omega$ are all ordinary vertices of the box. Every
weight is strictly positive, and every signed vertex value is
nonnegative by the induction hypothesis. If $f(x)=0$, every
vertex value must vanish. Interpolating again, now at any point of
the box, shows $f\equiv0$. If $d=0$, the box is a singleton and
the assertion is immediate.

For the rank assertion, all $(s+1)$-minors of the submatrix vanish
at the chosen interior point. Their order is at most $k-1$, so
each vanishes throughout the box. The rank is consequently at
most $s$ everywhere. If the submatrix has fewer than $s+1$ rows
or columns, the same conclusion is automatic.
\end{proof}

In particular, the rank at a relative-interior point is the maximum
rank of that submatrix over the box whenever this lemma applies.
The lemma uses all vertices only as an interpolation identity; it
does not add vertex hypotheses to Theorem~\ref{thm:main}.

\subsection{A three-term determinantal identity}

\begin{lemma}\label{lem:fekete}
Fix a $k$-element row set $R$ and a $(k-1)$-element subset
$R'\subset R$. Write a $k$-element column set as
$J=\{a\}\cup S\cup\{c\}$, where $a=\min J$, $c=\max J$,
and $|S|=k-2$. Let $a<b<c$ with $b\notin J$. For increasing
column sets define
\[
 \Delta(T)=\det Z[R,T],\qquad d(T)=\det Z[R',T]
\]
when the displayed determinants have the appropriate orders.
Then
\begin{align}
 \Delta(J)d(S\cup\{b\})
 ={}&\Delta((J\setminus\{a\})\cup\{b\})d(S\cup\{a\})\notag\\
 &+\Delta((J\setminus\{c\})\cup\{b\})d(S\cup\{c\}).
 \label{eq:fekete}
\end{align}
The identity has no rank hypothesis.
\end{lemma}

\begin{proof}
Let $T=J\cup\{b\}=\{t_1<\cdots<t_{k+1}\}$ and let $w_s$
be column $t_s$ restricted to the rows $R$. Laplace expansion of
a determinant with a repeated row gives the cofactor-kernel
identity
\[
 \sum_{s=1}^{k+1}(-1)^s\Delta(T\setminus\{t_s\})w_s=0.
\]
Restrict this vector identity to $R'$. Apply the linear functional
that appends its argument after the $k-2$ fixed columns indexed
by $S$ and takes the determinant. The terms indexed by $S$
vanish because a column is repeated. The other three terms are
those of $a,b,c$.

For $a$, moving the appended column to the first position crosses
$k-2$ columns; with its kernel sign the coefficient is
$(-1)^{k-1}$. If $b$ occupies position $t$ in $T$, then exactly
$k-t$ members of $S$ are larger than $b$. Its coefficient is
$(-1)^t(-1)^{k-t}=(-1)^k$. The column $c$ is already appended
in increasing order and has coefficient $(-1)^{k+1}$. Moving
the $b$ term to the other side gives~\eqref{eq:fekete}.
\end{proof}

If the $(k-1)$-minors have a common weak sign and
$d(S\cup\{b\})\neq0$, equation~\eqref{eq:fekete} expresses
$\Delta(J)$ as a nonnegative linear combination of two $k$-minors
with smaller column span. A vanishing denominator is the obstruction
that will produce a rank bound.

\subsection{Solid minors and minimal span}

\begin{lemma}\label{lem:solid}
Every solid $k$-minor has its prescribed sign throughout $\K$.
\end{lemma}

\begin{proof}
Let the solid row and column intervals start at $i_0,j_0$.
Use checker coordinates $z_{ij}=(-1)^{i+j}X_{ij}$, which
increase from $A$ to $B$. For an entry in local row $u$ and
column $v$ of this block, the derivative of its signed determinant
with respect to $z_{ij}$ is
\[
 \eps_k(-1)^{i+j+u+v}\det X[I\setminus\{i\},J\setminus\{j\}].
\]
Here $i=i_0+u-1$ and $j=j_0+v-1$, so the parity factor is the
constant $(-1)^{i_0+j_0}$. By the induction hypothesis, every
such derivative has weak sign
$(-1)^{i_0+j_0}\eps_k\eps_{k-1}$, independent of the chosen
entry. The signed determinant is thus monotone in one common
checker direction in all its free coordinates. Its minimum is
attained at $A$ or at $B$, where it is nonnegative.
\end{proof}

For a nonempty index set $I$, its span is $\max I-\min I$.

\begin{lemma}\label{lem:rectangle}
If a wrong-sign $k$-minor occurs in $\K$, it occurs at a matrix
$M\in\ri\K$. At that matrix choose a wrong-sign $k$-minor
minimizing the sum of its row and column spans. After transposition
if necessary, it can be written on rows $R$ and columns
\[
 J=\{a\}\cup S\cup\{c\},\qquad |R|=k,\quad |S|=r=k-2,
\]
with $C=[a+1,c-1]$ satisfying $|C|\geq r+1$ and
\begin{equation}\label{eq:rank-r}
 \rank M[R,C]=r.
\end{equation}
The columns $S$ are a basis of this rectangle, and
$\rank X[R,C]\leq r$ for every $X\in\K$.
\end{lemma}

\begin{proof}
A strictly wrong signed value persists under a sufficiently small
motion toward the midpoint of the box, so it occurs in $\ri\K$.
There are finitely many minors, so a minimum-span choice exists.
By Lemma~\ref{lem:solid}, its rows and columns cannot both be
consecutive. Transposition preserves the hypotheses and the sum
of the spans; we may arrange that $J$ has a gap. Then
$|C|\geq k-1=r+1$.

The chosen determinant is nonzero, so the columns $S$ restricted
to $R$ are independent. Suppose an omitted interior column
$b\in C\setminus S$ increased their rank to $r+1=k-1$.
Some $(k-1)$-element row set $R'\subset R$ would then give
$d(S\cup\{b\})\neq0$ in Lemma~\ref{lem:fekete}.
The two other $k$-minors in~\eqref{eq:fekete} have the same
rows and strictly smaller column span. They therefore have the
correct sign by minimality. All three $(k-1)$-minors have the
same weak sign, with a nonzero denominator, so the identity
would give the correct sign for the chosen $k$-minor as well.
This contradiction proves that every interior column is in the
span of $S$, establishing~\eqref{eq:rank-r}.

Since $r=k-2$, the permanent-zero lemma applies to all
$(r+1)$-minors in this rectangle. They vanish throughout the
box, proving the final rank assertion.
\end{proof}

\subsection{Independent entry variation and essential lines}

A row or column of a rank-$r$ matrix is \emph{essential} if
deleting it reduces the rank to $r-1$.

\begin{lemma}\label{lem:essential}
Let $F$ be a rectangular matrix of rank $r\geq1$.
If $\rank(F+tE_{ij})\leq r$ for all sufficiently small positive
and negative $t$, then row $i$ or column $j$ of $F$ is essential.
\end{lemma}

\begin{proof}
Every $r$-minor after deletion of row $i$ and column $j$ must
vanish. Otherwise, adjoining this row and column would give an
$(r+1)$-minor of $F+tE_{ij}$ whose coefficient of $t$ is, up to
sign, that nonzero $r$-minor. Its constant term is zero because
$\rank F=r$, so the rank bound would fail for nonzero small $t$.

Factor $F=UV$, where $U$ has $r$ independent columns and $V$
has $r$ independent rows. If neither indicated line were essential,
then deleting row $i$ from $U$ and column $j$ from $V$ would
leave both ranks equal to $r$. The first factor defines an
injective map from $\R^r$, and the second a surjective map onto
$\R^r$. Their product therefore has rank $r$. This is exactly
the submatrix of $F$ obtained by deleting both lines, in
contradiction to the vanishing of all its $r$-minors.
\end{proof}

Apply this lemma to the rectangle of Lemma~\ref{lem:rectangle}.
If one of its entries is free, the relative-interior matrix $M$
admits a small two-sided change in that entry alone. The permanent
rank bound makes Lemma~\ref{lem:essential} applicable. After
deleting the essential line, the residual rectangle has rank
$r-1$ at $M$. Its rank is at most $r-1$ throughout $\K$ by
Lemma~\ref{lem:permanent}, and in particular at both endpoints.

\section{Completion of the induction}\label{sec:induction}

\begin{proof}[Proof of Theorem~\ref{thm:main}]
Only (ii)$\Rightarrow$(i) needs proof. The case $n=1$ is immediate,
and Corollary~\ref{cor:nonsingular} proves nonsingularity throughout
the box for every $n$. Proposition~\ref{prop:ordertwo} supplies
the first two minor orders.

Suppose, inductively, that all minor orders through $k-1$ are
correct throughout $\K$, where $3\leq k\leq n$. Assume that a
wrong-sign $k$-minor exists and take $M,R,J,S,C,r$ from
Lemma~\ref{lem:rectangle}. Thus
\[
 J=\{a\}\cup S\cup\{c\},\qquad
 C=[a+1,c-1],\qquad r=k-2,
\]
and the rectangle $M[R,C]$ has rank $r$, with at least $r+1$
columns. Set
\[
 p=r-1=k-3,\qquad
 \delta=\eps_p\eps_{p+1}\eps_{p+2}\eps_{p+3}.
\]
We distinguish whether this interior-column rectangle is fixed
throughout the box.

\medskip\noindent
\textbf{Case 1: the rectangle is fixed.}
Choose an invertible $p\times p$ pivot $H=M[P,Q]$ with
$P\subset R$ and $Q\subset S$. Such a pivot exists because
$M[R,S]$ has rank $r>p$; when $p=0$, choose the empty pivot.
The pivot and every entry in $R\times C$ are fixed throughout
$\K$.

Using these same pivot positions, diagonal signs, and row reversal,
form the normalized Schur matrices
$\mathcal T_A,\mathcal T_B,\mathcal T_M$ of
Section~\ref{sec:schur}. They are nonsingular and
$\mathrm{TN}_2$. The endpoints have weak third-order sign
$\delta$. The bad $k$-minor becomes a three-minor of
$\mathcal T_M$ with the wrong sign relative to $\delta$.
Its columns are $a,q,c$, where $q$ is the unique element of
$S\setminus Q$, and its three rows are those of $R\setminus P$
in the order after any row reversal.

On these three rows and all columns of $C\setminus Q$, the
three Schur matrices share one fixed rank-one block. Every term
of the elimination formula for this block uses only entries in
$R\times C$, so it is fixed. The rank identity
\eqref{eq:rank-lifting} gives rank $r-p=1$.
Column $q$ is nonzero: applying the same rank identity to
$M[R,S]$ shows that its single residual column has rank one.
There are at least two columns in $C\setminus Q$, and the
surviving columns $a,c$ lie on opposite sides of all of them in
the full Schur column order.

We claim that the shared block contains at least two columns
positive on all three selected rows. If $\delta=-1$, no interior
column of either endpoint Schur matrix can contain a zero, by
Lemma~\ref{lem:twozeros}: its only possible zero entries occur
in the first and last columns. Thus every column of the block
is positive on these rows.

If $\delta=+1$, the bad three-minor of $\mathcal T_M$ is
negative. Lemma~\ref{lem:middle} forces its middle column $q$
to be positive on all three selected rows. Every other nonzero
column of the shared nonnegative rank-one block is a positive
multiple of this column. If no second nonzero column existed,
choose another column $b\in C\setminus Q$. It is zero on all
three selected rows. Staircase support then forces these rows
to be zero at $a$ when $b<q$, or at $c$ when $b>q$: in a row
support interval, a zero lying to one side of a positive entry
excludes every farther column on that side. Either conclusion
would make the target three-minor zero, a contradiction. This
proves the claim in both signs.

Choose column $q$ and a second positive column of the fixed
block. Together with the first and last of the three selected
rows, they form a positive vanishing two-minor at both endpoints.
Lemma~\ref{lem:corner} puts this anchor in a rank-one corner.
The orientation is the same at $A$ and $B$: its necessary index
inequalities exclude one of the two possible orientations for
both matrices. This corner contains all three selected rows
and reaches one column boundary. It therefore contains the
same outer column $d\in\{a,c\}$ at both endpoints, as well
as column $q$.

Lift the rank bound through the common pivot. The submatrix
on $R\times(S\cup\{d\})$ has just these three residual rows
and the residual columns $q,d$. It follows that
\begin{equation}\label{eq:fixed-rank}
 \rank A[R,S\cup\{d\}]\leq p+1=r,
 \qquad
 \rank B[R,S\cup\{d\}]\leq r.
\end{equation}
This endpoint statement must now be transferred to $M$.
Every $(r+1)$-minor in~\eqref{eq:fixed-rank} vanishes at both
endpoints. All its columns except possibly $d$ lie in the fixed
set $S$. Hence its determinant is affine in the entries of its
one remaining column $d$. It vanishes at the midpoint
$(A+B)/2$. Minors omitting $d$ are fixed and already zero.

The midpoint belongs to $\ri\K$, including when some entries
are fixed. The order $r+1=k-1$ is covered by the induction
hypothesis. Lemma~\ref{lem:permanent} therefore makes all these
midpoint zeros permanent throughout $\K$, and in particular
\[
 \rank M[R,S\cup\{d\}]\leq r.
\]
Adjoining the other outer column raises rank by at most one.
The purported bad $k$-minor consequently has rank at most
$r+1=k-1$, contradicting its nonzero determinant.
The affine midpoint step is essential here: endpoint zeros of
an arbitrary multiaffine determinant alone do not imply an
interior zero.

\medskip\noindent
\textbf{Case 2: the rectangle has a free entry and $\delta=-1$.}
Lemma~\ref{lem:essential} supplies an essential row or column.
After its deletion, the rank is at most $r-1=p$ at each
endpoint. Deleting a row leaves dimensions
\[
 (k-1)\times|C|\ \geq\ (p+2)\times(p+2),
\]
where the inequality compares both dimensions. Deleting a
column leaves dimensions
\[
 k\times(|C|-1)\ \geq\ (p+3)\times(p+1).
\]
Both violate Lemma~\ref{lem:rank-obstruction} at either
nonsingular sign-regular endpoint. This case is impossible,
even if the residual endpoint rank is strictly less than $p$.

\medskip\noindent
\textbf{Case 3: the rectangle has a free entry and $\delta=+1$.}
Again select an essential row or column. Schur elimination in
this case is performed at $M$ alone; its pivot is not required
to be fixed or to remain invertible at the endpoints.

Suppose first that row $i$ is essential. The columns $S$ span
the whole rectangle, so their restriction to $R\setminus\{i\}$
has rank $r-1=p$. Choose a nonsingular pivot on rows
$P\subset R\setminus\{i\}$ and columns $Q\subset S$ of
size $p$. On these remaining rows, every interior column lies
in the span of the pivot columns. Its Schur residual is
therefore zero outside the pivot rows. Among the three rows
of the target three-minor, the two other than $i$ are zero
in its middle column, namely the unique column in
$S\setminus Q$. The normalized Schur matrix is
$\mathrm{TN}_2$. Lemma~\ref{lem:middle} excludes the negative
determinant that a wrong sign relative to $\delta=+1$ would
require.

Suppose instead that column $j$ is essential. Every column
basis $S$ of the rectangle contains $j$, since deletion of
$j$ lowers rank. Put $Q=S\setminus\{j\}$ and choose $p$
rows $P\subset R$ giving an invertible pivot. The columns
$Q$ form a basis after deletion of $j$, so every other
interior column has zero Schur residual on the three
remaining selected rows. At least one such zero column
$b$ remains: $|C|\geq r+1=p+2$, while only $p$ pivot
columns and column $j$ have been excluded.

The target columns are $a,j,c$. A negative three-minor would
force column $j$ to be positive on all three selected rows by
Lemma~\ref{lem:middle}. The zero column $b$ and staircase
support then force column $a$ to be zero on these rows if
$b<j$, and column $c$ to be zero if $b>j$. The target
determinant is zero, a contradiction.

All three cases are excluded. Hence no wrong-sign $k$-minor
occurs in $\ri\K$. A strictly wrong value at a boundary point
would persist toward the midpoint, so none occurs on the
boundary either. This completes the induction. All prescribed
minor signs hold throughout $\K$, and its nonsingularity was
proved separately in Corollary~\ref{cor:nonsingular}.
\end{proof}

\section{Exact examples}\label{sec:examples}

The first three examples isolate the roles of nonsingularity,
the common signature, and the choice of vertices. The last gives a
four-parameter degenerate box outside the preceding fixed-entry
criteria. All entries and all checks in this section are rational.

\begin{example}[Singular endpoints]\label{ex:singular}
Consider
\[
 A=\begin{pmatrix}0&0&1\\0&0&0\\0&0&0\end{pmatrix},
 \qquad
 B=\begin{pmatrix}1&0&1\\0&0&0\\1&0&1\end{pmatrix}.
\]
Both are nonnegative of rank one, so every minor of order at
least two is zero and both matrices are totally nonnegative.
They satisfy $A\chle B$, since every differing position has
even index sum. At their midpoint $M=(A+B)/2$, however,
\[
 \det M[\{1,3\},\{1,3\}]
 =\det\begin{pmatrix}1/2&1\\1/2&1/2\end{pmatrix}
 =-\frac14.
\]
Thus singular checker vertices need not even preserve the weak
minor-sign conclusion. For every $n\geq3$, replacing $A,B$ by
$A\oplus I_{n-3},B\oplus I_{n-3}$ gives the same failure in
dimension $n$. A block-diagonal direct sum of totally nonnegative
matrices is totally nonnegative: each nonzero minor is a product
of minors in corresponding diagonal blocks, without an index
permutation. The displayed bad minor persists in the first block.
\end{example}

\begin{example}[A common signature is necessary]\label{ex:signature}
The matrices
\[
 A=\begin{pmatrix}1&2\\2&1\end{pmatrix},\qquad
 B=\begin{pmatrix}2&1\\1&2\end{pmatrix}
\]
are checker-comparable and strictly sign regular. Their signatures
are $(+1,-1)$ and $(+1,+1)$, respectively, because their entries
are positive and their determinants are $-3$ and $3$.
Their midpoint has every entry equal to $3/2$ and is singular.
This is an example in dimension two of the all-dimensional
theorem. Unlike total nonnegativity, arbitrary sign regularity
with a prescribed signature is not automatically preserved by
adjoining an identity block, so no such lifting is used here.
\end{example}

\begin{example}[Ordinary interval bounds]\label{ex:ordinary}
Let
\[
 L=\begin{pmatrix}1&2\\2&5\end{pmatrix},\qquad
 U=\begin{pmatrix}2&3\\3&6\end{pmatrix}.
\]
These are ordinary entrywise lower and upper bounds, and both
are strictly totally positive: their determinants are $1$ and $3$.
Nevertheless, the permitted matrix
\[
 X=\begin{pmatrix}1&3\\3&5\end{pmatrix}\in[L,U]
\]
has determinant $-4$. The ordinary lower and upper bounds
cannot replace the checker vertices in Theorem~\ref{thm:main}.
For $n\geq2$, the corresponding direct sums with $I_{n-2}$
give nonsingular totally nonnegative lower and upper bounds
with the same bad two-minor at an interior or boundary choice.
\end{example}

\begin{example}[A degenerate four-parameter box]\label{ex:degenerate}
Put
\[
 (\lambda_0,\lambda_1,\lambda_2,\lambda_3,\lambda_4)
 =\left(1,\frac1{10},\frac1{100},-\frac1{1000},-\frac1{10000}\right),
 \qquad
 q=\frac{500618646037}{19352067067828800},
\]
and define a rational $5\times5$ matrix by
\begin{equation}\label{eq:seed}
 W_{ij}=\sum_{s=0}^4\lambda_s\left(\frac{ij}{36}\right)^s
             -q\,\mathbf 1_{\{i=j=1\}},
 \qquad 1\leq i,j\leq5.
\end{equation}
This seed appears in the example in
\cite[Section~6]{HoldenGraph2026}. It is nonsingular and its
signature is
\[
 (+1,+1,+1,-1,+1).
\]
Its only zero minor is the
order-three minor on rows and columns $\{1,2,3\}$.

With $D_5=\diag(-1,1,-1,1,-1)$, let
\begin{equation}\label{eq:example-box}
 H=-D_5W^{-1}D_5,\qquad
 E=\begin{pmatrix}
 1&-1&0&0&0\\-1&1&0&0&0\\
 0&0&0&0&0\\0&0&0&0&0\\0&0&0&0&0
 \end{pmatrix},\qquad
 A=H,\quad B=H+\tfrac18 E.
\end{equation}
Then $A\chle B$. The four entries in the northwest $2\times2$
block vary independently, each over an interval of width $1/8$;
all other entries are fixed. Every entry of both endpoints is
strictly positive. Exact evaluation gives the counts in
Table~\ref{tab:example}, at \emph{each} endpoint. The unique
zero minor is
\begin{equation}\label{eq:zero-example}
 \det A[\{4,5\},\{4,5\}]
 =\det B[\{4,5\},\{4,5\}]=0.
\end{equation}
Thus both endpoints are nonsingular and have signature
\[
 \eps=(+1,+1,-1,+1,-1).
\]
Theorem~\ref{thm:main} applies to the whole four-parameter box.

\begin{table}[htbp]
\centering
\begin{tabular}{rrrr}
\toprule
Minor order & Positive & Zero & Negative\\
\midrule
1 & 25 & 0 & 0\\
2 & 99 & 1 & 0\\
3 & 0 & 0 & 100\\
4 & 25 & 0 & 0\\
5 & 0 & 0 & 1\\
\bottomrule
\end{tabular}
\caption{Exact minor counts at each checker vertex in
Example~\ref{ex:degenerate}. There are $251$ nonempty minors
in each $5\times5$ matrix.}\label{tab:example}
\end{table}

Here is why the earlier stated criteria do not cover this example.
The fixed positions $(3,3)$ and $(3,4)$ have different checker
parities, so the fixed-parity criterion does not apply.
The directed graph in~\cite{HoldenGraph2026} has an edge
$C_j\to R_i$ at a fixed positive entry of even index sum and
an edge $R_i\to C_j$ at one of odd index sum. It contains the
fixed-entry cycle
\[
 R_3\longrightarrow C_2\longrightarrow R_4
      \longrightarrow C_3\longrightarrow R_3.
\]
Thus Theorem~G's acyclicity hypothesis fails. The adjacent
$2\times2$ block in~\eqref{eq:zero-example} is fully fixed
and singular, so Theorem~V of~\cite{HoldenFive2026} also does
not apply. The endpoints are dense, and the vanishing two-minor
has four positive entries, so they are neither tridiagonal nor
almost strictly sign regular. Finally, their first four
signature signs $(+,+,-,+)$ do not occur in the eight
period-four families listed in~\cite{HassunehAdmGarloff2025};
changing the final sign of a five-dimensional signature does
not alter this exclusion.

The formulas~\eqref{eq:seed}--\eqref{eq:example-box} specify
the matrices exactly without rounded entries. The accompanying
standard-library rational verification program checks the seed,
both checker vertices, the zero pattern, and all sixteen
ordinary vertices of this box. At each of those sixteen
vertices, every signed minor is nonnegative and the signed
determinant is positive. Multiaffine interpolation therefore
also gives an independent finite verification of this
particular box. This finite calculation is an illustration;
the general theorem is established by the preceding proof.
\end{example}

\section{Scope and certification cost}\label{sec:scope}

Theorem~\ref{thm:main} concerns real square boxes with independent
scalar entry ranges and nonsingular endpoints of one prescribed
signature. It treats every dimension, every such signature, and
all allowable patterns of zero proper minors and fixed entries.
The rationality of the examples imposes no restriction on the
theorem. No perturbation that makes the proper minors strict
is required.

The result includes the usual two-vertex interval conclusion
for the nonsingular totally nonnegative, strictly sign-regular,
and special-signature classes discussed in the introduction.
It does not assert every stronger auxiliary conclusion of
those papers, such as preservation of a whole zero/nonzero
minor pattern. It also does not replace their results for
singular or rectangular matrices under additional hypotheses.
Example~\ref{ex:singular} shows that nonsingularity cannot be
dropped from the general two-vertex statement.

The equivalence gives a reduction in the number of matrices
whose nonsingular sign regularity must be certified: two
checker vertices suffice for the entire square box. Checking
that a given endpoint has the required signs is a separate
recognition problem. An arithmetic or bit-complexity claim
requires a specified endpoint-testing algorithm and a cost
analysis; a measured computational speedup requires an
implementation and comparison. Neither follows from the
vertex count alone. The present result is the exact
necessary-and-sufficient certification principle.

\section*{Acknowledgements}
[Acknowledgements to be decided]

\bibliographystyle{plainnat}
\bibliography{refs}
\end{document}
